% 写书 Skill 的 PDF 版式模板，由 scripts/build-pdf.py 作为 pandoc --include-in-header 传入。
% 下面的 BOOK-SETTINGS 一行会被替换为 typeset.yaml 生成的配色、封面与字体回退；
% 本书专用的补丁写进 typeset.yaml 的 pdf.extra_header，不要改这份模板。
\usepackage{fancyhdr}
\usepackage{eso-pic}
\usepackage{xcolor}
\usepackage{graphicx}
\usepackage{titlesec}
\usepackage{array}
\usepackage{ragged2e}
\usepackage{newunicodechar}

%% BOOK-SETTINGS %%

\newif\ifbookcoverpage
\bookcoverpagefalse

\fancypagestyle{fancy}{%
  \fancyhf{}
  \fancyfoot[R]{\color{bookfootergray}\fontsize{8}{10}\selectfont\thepage}
  \renewcommand{\headrulewidth}{0pt}
  \renewcommand{\footrulewidth}{0pt}
}
\fancypagestyle{plain}{%
  \fancyhf{}
  \fancyfoot[R]{\color{bookfootergray}\fontsize{8}{10}\selectfont\thepage}
  \renewcommand{\headrulewidth}{0pt}
  \renewcommand{\footrulewidth}{0pt}
}
\pagestyle{fancy}

% 扉页（\maketitle）按封面页处理
\let\bookoldmaketitle\maketitle
\renewcommand{\maketitle}{\bookcoverpagetrue\bookoldmaketitle\bookcoverpagefalse}

% 章标题在“：”处换行；没有冒号时在第一个全角空格处换行（如“前言　…”）；
% “附录 A　…”在全角空格处换行
\ExplSyntaxOn
\NewDocumentCommand\booksplittitle{m}
  {
    \tl_set:Nn \l_tmpa_tl {#1}
    \tl_if_in:NnT \l_tmpa_tl {附录}
      { \tl_replace_once:Nnn \l_tmpa_tl {　} {\\} }
    \tl_if_in:NnTF \l_tmpa_tl {：}
      { \tl_replace_once:Nnn \l_tmpa_tl {：} {：\\} }
      { \tl_if_in:NnF \l_tmpa_tl {附录} { \tl_replace_once:Nnn \l_tmpa_tl {　} {\\} } }
    \tl_use:N \l_tmpa_tl
  }
\ExplSyntaxOff
\titleformat{\chapter}[display]
  {\normalfont\huge\bfseries\color{bookink}\raggedright}
  {\color{bookaccent}\normalsize\mdseries\MakeUppercase{\chaptertitlename\ \thechapter}}
  {14pt}{\booksplittitle}
  [\vspace{6pt}{\color{bookaccent}\rule{72pt}{2pt}}]
\titlespacing*{\chapter}{0pt}{18pt}{40pt}

% 节标题超过一行时在“：”后换行，避免末行只剩一两个字
\ExplSyntaxOn
\NewDocumentCommand\booksplitsection{m}
  {
    \tl_set:Nn \l_tmpa_tl {#1}
    \hbox_set:Nn \l_tmpa_box {#1}
    \dim_compare:nNnT { \box_wd:N \l_tmpa_box } > { \linewidth }
      { \tl_replace_once:Nnn \l_tmpa_tl {：} {：\\} }
    \tl_use:N \l_tmpa_tl
  }
\ExplSyntaxOff
\titleformat{\section}
  {\normalfont\Large\bfseries\color{bookink}\raggedright}{\thesection}{0.8em}{\booksplitsection}
\titleformat{\subsection}
  {\normalfont\large\bfseries\color{bookink}\raggedright}{\thesubsection}{0.8em}{\booksplitsection}

% 图片默认居中且不超过版心
\setkeys{Gin}{width=\linewidth,height=0.9\textheight,keepaspectratio}
\makeatletter
\def\fps@figure{htbp}
\makeatother

% 表格行距
\renewcommand{\arraystretch}{1.28}

% 引用块样式
\usepackage{mdframed}
\renewenvironment{quote}
  {\begin{mdframed}[linewidth=0pt,leftline=true,linecolor=bookaccent,
     innerleftmargin=12pt,innerrightmargin=6pt,innertopmargin=8pt,
     innerbottommargin=8pt,backgroundcolor=black!3,skipabove=10pt,skipbelow=10pt]%
   \setlength{\parskip}{4pt}\small}
  {\end{mdframed}}

% 正文字体缺字的符号指定回退字体（字体名见 BOOK-SETTINGS）
\newunicodechar{☐}{{\booksymbolfont ☐}}
\newunicodechar{☑}{{\booksymbolfont ☑}}
\newunicodechar{☒}{{\booksymbolfont ☒}}
\newunicodechar{★}{{\booksymbolfont ★}}
\newunicodechar{○}{{\booksymbolfont ○}}
\newunicodechar{✓}{{\bookcheckfont ✓}}
\newunicodechar{◆}{{\bookcheckfont ◆}}
\newunicodechar{Ⅰ}{{\bookromanfont Ⅰ}}
\newunicodechar{Ⅱ}{{\bookromanfont Ⅱ}}
\newunicodechar{Ⅲ}{{\bookromanfont Ⅲ}}
\newunicodechar{Ⅳ}{{\bookromanfont Ⅳ}}
\newunicodechar{Ⅴ}{{\bookromanfont Ⅴ}}
\newunicodechar{①}{{\bookromanfont ①}}
\newunicodechar{②}{{\bookromanfont ②}}
\newunicodechar{③}{{\bookromanfont ③}}
\newunicodechar{④}{{\bookromanfont ④}}
\newunicodechar{⑤}{{\bookromanfont ⑤}}
\newunicodechar{⑥}{{\bookromanfont ⑥}}
\newunicodechar{⑦}{{\bookromanfont ⑦}}
\newunicodechar{⑧}{{\bookromanfont ⑧}}
\newunicodechar{⑨}{{\bookromanfont ⑨}}
\newunicodechar{⑩}{{\bookromanfont ⑩}}

% 封面必须是第一页，先于 \maketitle 输出
\ifdefined\bookcoverfile
\AtBeginDocument{%
  \bookcoverpagetrue
  \newgeometry{top=0mm,bottom=0mm,left=0mm,right=0mm}%
  \begin{titlepage}%
  \thispagestyle{empty}%
  \noindent\includegraphics[width=\paperwidth,height=\paperheight]{\bookcoverfile}%
  \end{titlepage}%
  \restoregeometry
  \bookcoverpagefalse
}
\fi

% 图题、表题已自带“图 X-Y”“表 X-Y”编号，去掉 LaTeX 自动前缀
\usepackage{caption}
\captionsetup{labelformat=empty,font={small,color=bookfootergray},skip=8pt}
\captionsetup[table]{position=top}

% 目录只列编与章；PDF 书签保留到节
\AtBeginDocument{\hypersetup{bookmarksdepth=2}}
% 目录中“编”的样式
\makeatletter
\renewcommand*\l@part[2]{%
  \addpenalty{-\@highpenalty}\vskip 1.0em \@plus\p@
  {\parindent\z@ \rightskip\@pnumwidth \parfillskip -\@pnumwidth
   \leavevmode\large\bfseries\color{bookaccent} #1\hfil\nobreak\hb@xt@\@pnumwidth{\hss #2}\par}
  \nobreak}
\makeatother

% 二三级标题前预留空间，避免标题孤立在页底
\usepackage{needspace}
\usepackage{etoolbox}
\makeatletter
% 紧跟在上一级标题之后的标题不再单独预留（否则会把上一级标题留在页底）
\preto\section{\if@nobreak\else\Needspace*{6\baselineskip}\fi}
\preto\subsection{\if@nobreak\else\Needspace*{5\baselineskip}\fi}
\@secpenalty=0
\makeatother
% 段落首末行不单独跨页
\clubpenalty=10000
\widowpenalty=10000
\displaywidowpenalty=10000
\usepackage{placeins}

% 浮动体：允许大图与正文同页；浮动页从顶端排起
\renewcommand{\topfraction}{0.92}
\renewcommand{\bottomfraction}{0.8}
\renewcommand{\textfraction}{0.06}
\renewcommand{\floatpagefraction}{0.88}
\setcounter{topnumber}{2}
\makeatletter
\setlength{\@fptop}{0pt}
\setlength{\@fpsep}{14pt}
\setlength{\@fpbot}{0pt plus 1fil}
\makeatother

% 网址可在任意位置断行，不插连字符
\usepackage{xurl}

% 中文弯引号：pandoc 输出的 ``…'' 已由 build-pdf.py 还原为 Unicode，这里让 xeCJK 按全角字形排
\xeCJKDeclareCharClass{FullLeft}{"2018, "201C}
\xeCJKDeclareCharClass{FullRight}{"2019, "201D}
\xeCJKsetwidth{“”‘’}{1em}
% 相邻引号：宋体的弯引号字形窄，默认挤压后会贴在一起，按字形留出可见间距
\xeCJKsetkern{”}{“}{0.3em}
\xeCJKsetkern{”}{‘}{0.3em}
\xeCJKsetkern{’}{“}{0.3em}
\xeCJKsetkern{’}{‘}{0.3em}
\xeCJKsetkern{’}{”}{0.12em}
\xeCJKsetkern{“}{‘}{0.12em}
\xeCJKsetkern{”}{？}{0.1em}
\xeCJKsetkern{”}{！}{0.1em}
% 西文连接号按西文处理
\xeCJKDeclareCharClass{Default}{"2013}
% 中文省略号用中文字体（居中）
\xeCJKDeclareCharClass{CJK}{"2026}
% 只含西文的引号（不经 xeCJK 全角处理）
\newcommand{\bookldq}{``}
\newcommand{\bookrdq}{''}
% “人工智能+”“互联网+”这类词的“+”：不与前字分开、不插中西文间距，书签里保留“+”
\newcommand{\bookplus}{\texorpdfstring{\nobreak\hbox{+}}{+}}

% 目录：为页码预留宽度，长条目自动折行；章条目续行与“第X章　”后的标题文字对齐
\makeatletter
\patchcmd{\l@chapter}{\setlength\@tempdima{1.5em}}{\setlength\@tempdima{4em}}{}{\PackageWarning{book}{l@chapter patch failed}}
\renewcommand{\@pnumwidth}{2.4em}
\renewcommand{\@tocrmarg}{3.4em}
\makeatother

% 不加载 babel（会接管中文标点），目录标题直接设为中文
\renewcommand{\contentsname}{目录}
